\scanpage{64}
% archive sheet number: 63

soit $X'=\hat X-\mathring T$. qui est cpct. \deleted{\unsure{et $\in C$.|$X\in\mathcal C$.}} On vérifie

\margin{\uline{vérifier}}

\[
\begin{aligned}
\text{a)}\quad & X'\in C_\infty \qquad\text{donc}\quad X'\in \unsure{\mathcal E|\mathscr E}_\chi\\
\text{b)}\quad & X'\sim X=\hat X-\omega \qquad\text{donc}\quad X\in\mathcal E_\chi,
\end{aligned}
\]

\uline{NB} On trouve

\[
\chi(X)=\chi(\hat X-\mathring T)=\chi(\hat X)-\underset{\deleted{\unsure{\Vert\,1}}}{\underline{\chi(T)}}+\chi(\partial T)
\]

et $\chi(T)=1$ car $T$ \unsure{contractible|contr.ible|contractile}, donc

\[
\chi(X)=\chi(\hat X)-1+\chi(\underset{\substack{\Vert\\ S_\omega}}{\underline{\partial T}})=\chi_!(X)+\chi(\partial T)\qquad\text{i.e}
\]

\diagram{\fbox{\begin{tikzcd}[ampersand replacement=\&,column sep=0pt,row sep=1ex]
\chi(X)-\chi_!(X)=\chi( \& S_\omega\arrow[dr,bend right=38] \& )\\
\& \& \substack{\text{exprime la ``structure}\\ \text{locale de $X$ à l'infini''}}
\end{tikzcd}}}
